\documentclass[11pt]{article}

\usepackage{amsmath,amssymb,amsthm,mathtools}
\usepackage[T1]{fontenc}
\usepackage{lmodern}
\usepackage[margin=1in]{geometry}
\usepackage{enumitem}
\usepackage{hyperref}
\usepackage{microtype}

\hypersetup{
  colorlinks=true,
  linkcolor=blue,
  citecolor=blue,
  urlcolor=blue
}

\newtheorem{theorem}{Theorem}[section]
\newtheorem{proposition}[theorem]{Proposition}
\newtheorem{corollary}[theorem]{Corollary}
\newtheorem{definition}[theorem]{Definition}
\newtheorem{remark}[theorem]{Remark}

\newcommand{\Z}{\mathbb Z}

\title{A cleaned theorem suite in independent form for odd-$m$ D5}
\author{}
\date{March 2026}

\begin{document}
\maketitle

\begin{abstract}
This note repackages the current odd-$m$ D5 closure into one
manuscript-order suite of theorem statements. The statements are no longer
split across the
\[
033 \to 062 \to 076 \to 077 \to 079 \to 081 \to 082 \to 083
\]
chain. The structural first-exit theorem and the chart-to-raw composition step
are written explicitly, and the final gluing argument is presented as one
continuous proof.

The note stays honest about support. The componentwise concrete bridge,
fixed-$\delta$ tail-length reduction, chart/interface theorem, and regular
continuation theorem are retained in their accepted theorem form rather than
rederived here from the original compute-backed files.
\end{abstract}

\section{Scope and conventions}

Fix an odd modulus $m$ in the accepted D5 regime, and fix the best-seed channel
\[
R1 \to H_{L1}
\]
with endpoint seed
\[
L=[2,2,1], \qquad R=[1,4,4].
\]

The final theorem we want in cleaned form is:
\begin{quote}
Every actual lift of the exceptional cutoff row $\delta=3m-3$ continues through
$3m-2 \to 3m-1$ and lies in the regular continuing endpoint class there.
\end{quote}
Equivalent corollaries are:
\begin{enumerate}[label=\textup{(\arabic*)},leftmargin=2.5em]
\item no mixed-status realized $\delta$ exists on the true accessible union;
\item fixed realized $\delta$ determines tail length;
\item fixed realized $\delta$ determines endpoint class and padded future word;
\item $\rho = \rho(\delta)$ globally;
\item raw global $(\beta,\delta)$ is exact.
\end{enumerate}

The goal of this note is not to enlarge the accepted scope. It is to put the
current odd-$m$ D5 result into a theorem package that is independent in
statement and dependency order.

\section{Definitions}

\begin{definition}[full chains and actual lifts]
A \emph{full chain} is a complete boundary block between consecutive visits to
the $\Theta=2$ section. Its boundary label is the intrinsic digit $\delta$
defined below.

An \emph{actual lift} of a row $\delta_0$ is an actual accessible full-chain
occurrence whose boundary label is $\delta_0$.
\end{definition}

\begin{definition}[accessible boundary union and components]
The \emph{accessible boundary union} is the union of all actual accessible full
chains.

A \emph{splice-connected accessible component} is a connected component of that
union under the successor/predecessor splice relation between consecutive full
chains.
\end{definition}

\begin{definition}[tail length]
The \emph{remaining tail length} of an actual lift is the number of full-chain
successors still available before terminality. If successors continue forever,
its tail length is $\infty$.
\end{definition}

\begin{definition}[endpoint class and $\rho$]
Choose the standard accepted terminal padding convention on future current-event
words, so that finite and infinite tails are compared in one common one-sided
shift space.

Two actual lifts are in the same \emph{endpoint class} when they determine the
same padded future current-event word. The corresponding boundary
right-congruence class is written $\rho$.
\end{definition}

\begin{definition}[distinguished rows]
The \emph{exceptional cutoff row} is
\[
\delta = 3m-3.
\]
The corresponding interface rows are
\[
3m-2 \qquad \text{and} \qquad 3m-1.
\]
\end{definition}

\section{Structural first-exit package}

\begin{theorem}[structural first-exit theorem]
Let $\rho = u_{\mathrm{source}} + 1 \pmod m$ for the chosen source family, and
let $\widehat x_n=(q_n,w_n,u_n,\Theta_n)$ be the candidate active orbit
starting from the alt-$2$ entry
\[
E=(m-1,1,u_{\mathrm{source}},2).
\]
Then the following hold.
\begin{enumerate}[label=\textup{(\arabic*)},leftmargin=2.5em]
\item The trigger family is exactly
\[
H_{L1}=\{(q,w,u,\lambda)=(m-1,m-1,u,2):u\neq 2\}.
\]
\item Every phase-$1$ state on the candidate orbit satisfies
\[
q \equiv u-\rho+1 \pmod m.
\]
\item The only phase-$1$ states from which an alternate step can land in
$H_{L1}$ are
\[
T_{\mathrm{reg}}=(m-1,m-2,1) \quad \text{by direction } 2,
\]
and
\[
T_{\mathrm{exc}}=(m-2,m-1,1) \quad \text{by direction } 1.
\]
\item If $\rho\neq 4$, the first exit is $T_{\mathrm{reg}}$. If $\rho=4$, the
first exit is $T_{\mathrm{exc}}$.
\item The actual active branch agrees with the candidate orbit up to first
exit. Therefore the regular families first exit at $T_{\mathrm{reg}}$ by
direction $2$, and the exceptional family first exits at $T_{\mathrm{exc}}$ by
direction $1$.
\end{enumerate}
\end{theorem}

\begin{proof}
Write the candidate update as
\[
\widehat F(q,w,u,0)=(q+1,w,u,1),
\]
\[
\widehat F(q,w,u,1)=(q,w,u+1,2),
\]
\[
\widehat F(q,w,u,2)=\bigl(q,w+\mathbf 1_{\{q=m-1\}},u,3\bigr),
\]
\[
\widehat F(q,w,u,\Theta)=(q,w,u,\Theta+1)
\qquad (3\le \Theta\le m-1).
\]

The exact trigger family is the promoted $033$ formula, namely
\[
H_{L1}=\{(m-1,m-1,u,2):u\neq 2\}.
\]
So only phase-$1$ current states can trigger an alternate step into $H_{L1}$,
and a trigger must arrive at $(q',w')=(m-1,m-1)$.

On the candidate orbit, every phase-$1$ state satisfies
\[
q \equiv u-\rho+1 \pmod m.
\]
This is the phase-$1$ source-residue invariant of $062$: it holds at the first
phase-$1$ state after entry, and from one phase-$1$ state to the next both
sides increment by $1$.

Now there are only two ways to reach $(q',w')=(m-1,m-1)$ at layer $2$:
\begin{itemize}[leftmargin=2em]
\item by direction $2$ from $(q,w,\Theta)=(m-1,m-2,1)$,
\item by direction $1$ from $(q,w,\Theta)=(m-2,m-1,1)$.
\end{itemize}
These are exactly $T_{\mathrm{reg}}$ and $T_{\mathrm{exc}}$. At
$T_{\mathrm{reg}}$ the invariant gives
\[
u \equiv \rho-2 \pmod m,
\]
so the trigger condition $u\neq 2$ holds exactly when $\rho\neq 4$. At
$T_{\mathrm{exc}}$ the invariant gives
\[
u \equiv \rho-3 \pmod m,
\]
so when $\rho=4$ one has $u\equiv 1$, hence $u\neq 2$, and the exit occurs
there. Thus the candidate first exits are the universal targets stated above.

The remaining $062$ input is pre-exit patch avoidance: before the candidate
first-exit time, the candidate orbit avoids the patched current classes. Hence
all pre-exit actual states are $B$-labeled and therefore follow the same
full-coordinate updates as the candidate orbit. So the actual active branch
agrees with the candidate orbit up to first exit, and the same first-exit
statement holds on the actual branch.
\end{proof}

\section{Accepted support theorems in cleaned form}

\begin{proposition}[componentwise concrete bridge; accepted support theorem]
On the unique $\Theta=2$ state $(q,w,u,2)$ of a full regular chain, define
\[
\sigma \equiv w+u-q-1 \pmod m,
\qquad
\delta := q+m\sigma.
\]
Then on each splice-connected accessible component:
\begin{enumerate}[label=\textup{(\arabic*)},leftmargin=2.5em]
\item consecutive full chains satisfy the odometer splice law
\[
\delta' = \delta+1 \pmod{m^2};
\]
\item the current event class $\epsilon_4$ is determined by $(\beta,\delta)$;
\item the realized boundary image is a forward interval in $\Z/m^2\Z$ or the
full odometer cycle.
\end{enumerate}
\end{proposition}

\begin{proof}[Proof status]
This is the promoted theorem content of $076$, rewritten as one compact bridge
statement.
\end{proof}

\begin{proposition}[tail-length reduction; accepted support theorem]
At fixed realized $\delta$, two actual lifts can differ only by remaining
full-chain tail length.

Equivalently, once the concrete bridge package is accepted, there is no further
local future-law ambiguity at fixed realized $\delta$; the only possible
ambiguity is endpoint geometry.
\end{proposition}

\begin{proof}[Proof status]
This is the promoted theorem content of $077$.
\end{proof}

\begin{proposition}[chart/interface theorem; accepted support theorem]
The exceptional continuation is pinned in chart / chain-label coordinates to
\[
3 \to (\text{terminal chain of }4)\to 1,
\]
equivalently in $\delta$ coordinates to the interface
\[
3m-3 \to 3m-2 \to 3m-1.
\]
\end{proposition}

\begin{proof}[Proof status]
This is the promoted theorem content of $079$.
\end{proof}

\begin{proposition}[regular continuation theorem; accepted support theorem]
Every regular actual lift of every realized $\delta$ continues. In particular,
regular rows cannot create a second endpoint sheet at fixed $\delta$.
\end{proposition}

\begin{proof}[Proof status]
This is the promoted theorem content of $081$.
\end{proof}

\section{Chart-to-raw composition}

\begin{theorem}[actual exceptional continuation theorem]
Assume the structural first-exit theorem and the accepted chart/interface
theorem. Then every actual lift of the exceptional cutoff row
\[
\delta=3m-3
\]
continues through the actual interface rows
\[
3m-2 \to 3m-1.
\]
\end{theorem}

\begin{proof}
By the structural first-exit theorem, the exceptional source family is exactly
the family $\rho=4$, and its first exit occurs at the universal exceptional
target $T_{\mathrm{exc}}$ by direction $1$.

The accepted chart/interface theorem identifies the post-exit continuation of
that family in chain-label coordinates as
\[
3 \to (\text{terminal chain of }4)\to 1,
\]
equivalently in $\delta$ coordinates as
\[
3m-3 \to 3m-2 \to 3m-1.
\]

Because the structural theorem already identifies the exceptional family on the
actual branch and the chart theorem identifies the unique continuation of that
family after first exit, the chart/interface landing composes directly with the
raw first-exit theorem. Hence every actual exceptional lift of $3m-3$ continues
through $3m-2$ and lands in the regular continuing class at $3m-1$.
\end{proof}

\section{Final gluing theorem}

\begin{theorem}[final gluing theorem in cleaned form]
On the true accessible boundary union of odd-$m$ D5, every realized
$\delta$ has realization-independent padded future current-event word.
Equivalently:
\begin{enumerate}[label=\textup{(\arabic*)},leftmargin=2.5em]
\item no mixed-status realized $\delta$ exists;
\item fixed realized $\delta$ determines tail length;
\item fixed realized $\delta$ determines endpoint class;
\item $\rho=\rho(\delta)$ globally;
\item raw global $(\beta,\delta)$ is exact.
\end{enumerate}
\end{theorem}

\begin{proof}
By the regular continuation theorem, every regular actual lift continues. So no
regular row can create mixed status at fixed $\delta$.

The only remaining possible mixed-status row is the exceptional cutoff row
\[
\delta=3m-3.
\]
But by the actual exceptional continuation theorem, every actual lift of
$3m-3$ continues through the interface
\[
3m-2 \to 3m-1
\]
and lands in the regular continuing class there. Therefore the exceptional row
also does not create a second endpoint sheet.

Hence no realized $\delta$ has mixed continuation status on the true accessible
union. By the accepted tail-length reduction theorem, fixed realized $\delta$
can differ only by tail length; therefore no mixed status implies
realization-independent tail length. With the standard accepted terminal padding
convention, that implies realization-independent endpoint class and padded
future current-event word. So $\rho=\rho(\delta)$ globally.

Finally, the concrete bridge proposition says that on each component the current
event class is already determined by $(\beta,\delta)$ and the boundary dynamics
follow the odometer splice law. Since $\rho$ is now globalized as a function of
$\delta$ alone, the componentwise concrete bridge upgrades to raw global
$(\beta,\delta)$.
\end{proof}

\begin{corollary}[exact verifiable odd-$m$ solution]
For odd $m$ in the accepted D5 regime, raw global $(\beta,\delta)$ is exact on
the true accessible boundary union.
\end{corollary}

\begin{proof}
This is the final assertion of the gluing theorem.
\end{proof}

\section{What is still imported rather than reproved}

\begin{remark}
This note is independent in theorem order, but not yet completely raw from the
earliest files. The following inputs are still imported in accepted-support
form:
\begin{enumerate}[label=\textup{(\arabic*)},leftmargin=2.5em]
\item the componentwise concrete bridge package ($076$),
\item the fixed-$\delta$ tail-length reduction ($077$),
\item the chart/interface theorem ($079$),
\item the regular continuation theorem ($081$),
\item the patch-avoidance clause inside the structural first-exit theorem.
\end{enumerate}
So this is the first cleaned independent theorem suite, but not yet the final
self-contained replacement for every support theorem.
\end{remark}

\end{document}
